\section{A source-faithful rational degree-one tripod realization}
\label{sec:iut-rational-degree-one-source-realization}

The shadow construction of
Section~\ref{sec:iut-rational-tripod-shadow} can be compared directly with the
source definitions in Mochizuki's \emph{Arithmetic Elliptic Curves in General
Position}.  Put
\[
 P=\mathbb P^1_{\Q},\qquad C=\{0,1,\infty\},\qquad U=P\setminus C,
 \qquad U_{(0,1)}(\Q)=\{x\in\Q:0<x<1\}.
\]
The map $j(x)=[x:1]$ is an injective map to $U(\Q)$.  Its points have minimal
field of definition $\Q$, and hence exact degree one.  The recurrence in the
abstract LANA height theory already forces
\begin{equation}\label{eq:ptle-one-equals-pteq-one}
                    \operatorname{ptLE}(X,1)=\operatorname{ptEQ}(X,1),
\end{equation}
because the degree-at-most-zero locus is empty.

Let $h$ be the normalized logarithmic Weil height and let $n$ be the truncated
count for the three marked points.  For the genuine arithmetic objects, one
has
\begin{equation}\label{eq:rational-source-normalizations}
 \operatorname{ht}_{\omega_P(C)}\circ j\approx h,
 \qquad \log\operatorname{-diff}_P\circ j=0,
 \qquad \log\operatorname{-cond}_C\circ j\approx n.
\end{equation}
Indeed, $\omega_P(C)\cong\mathcal O_{\mathbb P^1}(1)$, and invariance of the
height BD class under isomorphism gives the first comparison.  The different
ideal of $\Q/\Q$ is the unit ideal, giving the middle equality.  Finally, if
$x=m/r$ is reduced with $0<m<r$, pulling $0,1,\infty$ back to
$\operatorname{Spec}\mathbb Z$ gives the divisors of $m,r-m,r$.  Reduction of
the pullback divisor therefore gives
\[
 \log\operatorname{-cond}_C(j(x))
   =\sum_{p\mid m(r-m)r}\log p=n(x)
\]
for the standard model.  Changing the proper flat model changes this function
only by bounded discrepancy.

For an abstract height theory $T$, call a map
$j:U_{(0,1)}(\Q)\to T.\operatorname{Pt}(T.\operatorname{tripod})$ a
source-faithful degree-one realization when it has exact-degree-one image and
satisfies the three relations in \eqref{eq:rational-source-normalizations}
after replacing the arithmetic functions by the corresponding fields of $T$.
Write $H_T,D_T,N_T$ for the three pulled-back functions.  Define the restricted
Statement~I by
\begin{equation}\label{eq:rational-source-restricted-statement-i}
 \forall\epsilon>0\ \exists C_\epsilon\ \forall x,\qquad
 H_T(x)\le(1+\epsilon)(D_T(x)+N_T(x))+C_\epsilon.
\end{equation}
We denote this proposition by $\operatorname{RStmtI}(T,j)$.

\begin{theorem}[Source-faithful degree-one transport]
\label{thm:rational-source-restricted-iff-abc}
For every source-faithful degree-one realization,
\[
 \operatorname{RStmtI}(T,j)
 \quad\Longleftrightarrow\quad \mathrm{ABCConjecture}.
\]
Consequently, full $T.\mathrm{StatementI}$ together with such a realization
implies the integer $abc$ conjecture.
\end{theorem}

\begin{proof}
Choose uniform constants $A_+,A_-,B_+,B_-$ such that
\[
 H_T\le h+A_+,\quad h\le H_T+A_-,\qquad
 N_T\le n+B_+,\quad n\le N_T+B_-.
\]
If \eqref{eq:rational-source-restricted-statement-i} holds, then, using
$D_T=0$,
\[
 h(x)\le(1+\epsilon)n(x)+
 C_\epsilon+A_-+(1+\epsilon)B_+.
\]
This is the uniform rational-tripod formulation of integer $abc$.  Conversely,
that formulation gives
\[
 H_T(x)\le(1+\epsilon)N_T(x)+
 C_\epsilon+A_++(1+\epsilon)B_-,
\]
which is \eqref{eq:rational-source-restricted-statement-i}.  Finally, full
Statement~I restricts along $j$ because
\eqref{eq:ptle-one-equals-pteq-one} puts every image point in its degree-one
domain.
\end{proof}

This theorem closes the elementary source normalization but not the IUT
argument.  At the pinned versions, \texttt{Genl.HeightTheory} is only an
abstract record, no genuine arithmetic instance is constructed, and the pinned
IUT repository has no connection to Statement~I.  A later public IUT snapshot
has only a conditional connection that still takes the height theory, its proof
package, and the IUT and analytic inputs as premises.

The distinction cannot be erased by adding the abstract proof package.  The
degree-empty, height-zero, and different-inflated pressure models all admit
full \texttt{ProofPackage} instances.  They show respectively
that this package does not force rational degree-one points, the canonical
height normalization, or the rational different.  A separate
conductor-inflated model preserves exact points, degree, height, and zero
different while destroying every uniform upper conductor comparison.  These
are full-premise countermodels to their precisely stated weakened extraction
claims; they do not contradict the genuine computations above, IUT, or $abc$.
For clarity, the first three models retain the abstract proof package.  The
conductor-inflated model retains the height theory, Statement~I, and the stated
point, degree, height, and different data, but no proof-package instance is
claimed for it.

The Lean companion proves the recurrence
\eqref{eq:ptle-one-equals-pteq-one}, formalizes the abstract transport theorem,
constructs the rational-shadow consistency witness, and checks the pressure
models without placeholders or new axioms.  Constructing the genuine
arithmetic height-theory instance and proving its Statement~I remain active
upstream tasks.
